\chapter{Biblioteca de Padrões, Checklists e Mandamentos}
\label{ap:padroes}

\section{Os 10 mandamentos da Engenharia de Prompts}

Uma versão condensada da filosofia deste livro:

\begin{enumerate}[leftmargin=*]
\item \textbf{Não terás outros oráculos além dos dados}
LLMs são modelos estatísticos, não fontes mágicas de verdade. Sempre que o domínio for crítico, valida com dados e especialistas.
\item \textbf{Não farás prompts vagos quando quiseres resultados específicos}
Define objetivo, formato, constraints. “Faz qualquer coisa fixe” é convite a lixo.
\item \textbf{Honrarás o contexto e o exemplo}
Bons exemplos e contexto relevante valem mais do que duas páginas de teoria no system prompt.
\item \textbf{Não tomarás a “primeira resposta” como resposta definitiva}
Itera. Pede revisão, pede alternativas, pede auto-crítica do modelo.
\item \textbf{Lembrar-te-ás de que modelos mudam}
O que funciona hoje pode mudar amanhã com uma nova versão. Por isso, testa, versiona e loga.
\item \textbf{Guardarás os teus prompts como código}
Em repositórios, com histórico, owners, testes. Não em caixas de texto anónimas numa UI qualquer.
\item \textbf{Terás guardrails fora do modelo}
Filtros, validações, controlos de acesso. System prompt forte ajuda, mas não chega.
\item \textbf{Usarás IA para pensar melhor, não para abdicar de pensar}
Brainstorm e ajuda na escrita, sim. Delegar julgamento ético e responsabilidade, não.
\item \textbf{Respeitarás quem é afetado pelos teus sistemas}
Utilizadores, clientes, pessoas cujos dados passam pelo teu pipeline. Transparência, fundamento de licitude e mitigação de danos.
\item \textbf{Partilharás o que aprendes}
Documenta padrões, anti-padrões, lições aprendidas. Engenharia de prompts é, em grande parte, conhecimento tácito transformado em boas práticas.
\end{enumerate}

\section{Padrões essenciais por tipo de tarefa}

Não é uma lista exaustiva, mas uma mini “carta de bolso”.

\subsection*{Design e clarificação de tarefa}

\begin{itemize}[leftmargin=*]
\item \textbf{Clarificar objetivo}:
\begin{lstlisting}
Vou descrever o que quero fazer. Primeiro,
faz apenas perguntas para clarificar o objetivo,
escopo e constraints. Não dês ainda solução.

[descrição livre]
\end{lstlisting}
\item \textbf{Reformular problema}:
\begin{lstlisting}
Reformula o problema abaixo em 2-3 versões:

Versão simples (para não especialistas).

Versão técnica (para equipa de engenharia).

Versão operacional (para equipa de operações).

[problema]
\end{lstlisting}
\end{itemize}

\subsection*{Geração de texto estruturado}

\begin{itemize}[leftmargin=*]
\item \textbf{Formato rígido}:
\begin{lstlisting}
Responde EXACTAMENTE no formato JSON abaixo,
sem texto adicional:

{
"categoria": "...",
"prioridade": "...",
"precisa_humano": false,
"resumo": "..."
}

Texto:
[text]
\end{lstlisting}
\item \textbf{Checklist}:
\begin{lstlisting}
A partir do texto abaixo, gera uma checklist
operacional, em bullet points, que alguém
poderia seguir passo a passo.

[texto]
\end{lstlisting}
\end{itemize}

\subsection*{Código e engenharia de software}

\begin{itemize}[leftmargin=*]
\item \textbf{Especificar antes de gerar}:
\begin{lstlisting}
Antes de escreveres qualquer código, quero que:

listes os requisitos funcionais;

listes os requisitos não funcionais (desempenho,
segurança mínima, logging);

proponhas uma estrutura de ficheiros e módulos.

Só depois de eu dizer "ok", gera o código.
\end{lstlisting}
\item \textbf{Explanação de código legado}:
\begin{lstlisting}
Explica o seguinte ficheiro de código:

o que faz em alto nível;

quais as funções principais e seus argumentos;

quaisquer riscos óbvios (segurança, performance),
se existirem.

[código]
\end{lstlisting}
\end{itemize}

\subsection*{Dados, ETL e análise}

\begin{itemize}[leftmargin=*]
\item \textbf{Schema-first}:
\begin{lstlisting}
O teu objetivo é extrair dados em JSON com o
seguinte exemplo de instancia; os tipos e as opcoes
devem constar de um esquema separado:

{
"nome": "Ana",
"idade": null,
"cidade": "Lisboa"
}

Se um campo não estiver presente no texto,
usa null.

Texto:
[text]
\end{lstlisting}
\item \textbf{EDA guiada}:
\begin{lstlisting}
Vou colar abaixo um resumo estatístico e
alguns gráficos (em texto). Age como analista
de dados.

Resume em 5-7 bullets o que os dados parecem
mostrar.

Propõe 3 hipóteses plausíveis.

Sugere 3 análises adicionais a fazer com os
dados originais.

[output da ferramenta]
\end{lstlisting}
\end{itemize}

\subsection*{Educação e explicação}

\begin{itemize}[leftmargin=*]
\item \textbf{Níveis de explicação}:
\begin{lstlisting}
Explica o conceito abaixo em três níveis:

Para um aluno do secundário.

Para um estudante de licenciatura na área.

Para um profissional que precisa de refrescar
o conceito rapidamente.

[conceito]
\end{lstlisting}
\item \textbf{Exercícios com soluções}:
\begin{lstlisting}
Gera 5 exercícios de dificuldade crescente
sobre [tema]. Para cada um:

enunciado,

solução detalhada,

2-3 erros comuns a evitar.

Formata em português de Portugal.
\end{lstlisting}
\end{itemize}

\subsection*{Negócios, gestão e comunicação profissional}

\begin{itemize}[leftmargin=*]
\item \textbf{E-mail em tom ajustado}:
\begin{lstlisting}
A partir dos pontos soltos abaixo, escreve
um e-mail em português de Portugal para
[destinatário], com tom:

[formal/profissional/relaxado mas respeitoso]

Inclui:

contexto breve,

pedido claro,

próximos passos.

[Pontos soltos]
\end{lstlisting}
\item \textbf{Análise crítica}:
\begin{lstlisting}
Lê o texto estratégico abaixo e responde:

3 pontos fortes.

3 pontos fracos/risco.

3 perguntas que faltam responder para tornar
o plano mais concreto.

[texto]
\end{lstlisting}
\end{itemize}

\subsection*{Segurança, políticas e compliance}

\begin{itemize}[leftmargin=*]
\item \textbf{Resumir políticas}:
\begin{lstlisting}
A partir desta política interna:

resume em 5-10 bullets as obrigações práticas
para um colaborador normal;

destaca 3 comportamentos proibidos;

sugere 3 exemplos concretos de "isto pode" e
3 de "isto não pode".

[política]
\end{lstlisting}
\item \textbf{Simular auditor}:
\begin{lstlisting}
Assume o papel de auditor de conformidade.
Lê a descrição do sistema abaixo e indica:

dados pessoais envolvidos,

potenciais riscos para privacidade,

5 perguntas que terias para a equipa.

[descrição do sistema]
\end{lstlisting}
\end{itemize}

\section{Checklist rápida de design de prompt}

Uma versão “cola na parede”:

\begin{itemize}[leftmargin=*]
\item \textbf{1. Objetivo}
Sei exatamente o que quero da resposta? (tipo de tarefa + público-alvo)
\item \textbf{2. Contexto}
Dei contexto suficiente? (exemplos, constraints, dados relevantes)
\item \textbf{3. Formato}
Especifiquei claramente o formato de saída? (JSON, bullets, secções)
\item \textbf{4. Limites}
Disse o que o modelo \emph{não} deve fazer? (tabus, escopo)
\item \textbf{5. Exemplos}
Incluí 1–3 exemplos representativos (se útil)?
\item \textbf{6. Robustez}
Pensei em inputs estranhos? (erros, abreviações, casos de canto)
\item \textbf{7. Segurança}
Indiquei restrições de conteúdo/sensibilidade quando relevante?
\item \textbf{8. Iteração}
Fiz pelo menos uma ronda de teste+ajuste com casos reais?
\end{itemize}

\section{Checklist rápida de LLMOps para um fluxo em produção}

Para cada fluxo baseado em LLM no teu sistema:

\begin{itemize}[leftmargin=*]
\item \textbf{Prompts versionados}
Estão em repositório, com histórico e owners?
\item \textbf{Golden set}
Existe um conjunto de teste, ainda que pequeno?
\item \textbf{Métricas}
Sei o que estou a otimizar? (accuracy, custo, satisfação, etc.)
\item \textbf{Logs e privacidade}
Tenho logs suficientes para debug, sem violar privacidade?
\item \textbf{Guardrails}
Há filtros de input/output e limites claros para ferramentas?
\item \textbf{Rollout}
Mudo de prompt/modelo com testes prévios e canary release?
\item \textbf{Incidentes}
Há plano para lidar com respostas perigosas ou erros graves?
\end{itemize}
